
\question
The protein that is activated in response to auxin that allows phototropism to occur is

\begin{choices}
\choice proton channel. 
\choice basal apparatus. 
\choice abscisic acid. 
\choice transmembrane. 
\correctchoice expansin. 
\choice wallulase. 
\end{choices}

\question
Which of the following about phloem is false?

\begin{choices}
\choice The cell has cytoplasm. 
\choice Sugars move by passive diffusion. 
\choice Sieve-tube elements do not have a nucleus. 
\correctchoice The plasma membrane is modified to form a sieve plate. 
\choice Companion cells provide proteins to sieve-tube elements. 
\end{choices}

\question
During protist conjugation, which of the following does \emph{not} occur during some part of the process?

\begin{choices}
\choice Cells exchange micronuclei.  
\correctchoice The macronucleus fuses with the micronucleus. 
\choice The micronucleus undergoes meiosis. 
\choice The micronucleus undergoes mitosis. 
\choice Some micronuclei distintegrate. 
\choice The macronucleus distintegrates. 
\end{choices}

\question
When ATP attaches to the myosin head

\begin{choices}
\choice the myosin head attaches to the actin filament.  
\choice the actin head attaches to the myosin filament. 
\correctchoice the myosin head releases from the actin filament. 
\choice the actin head attaches to the myosin filament. 
\choice the myosin head initiates the power stroke. 
\end{choices}

\question
Which of the following are primarily responsible for movement in the eukaryotic flagellum?

\begin{choices}
\choice Hydrogen ions and basal apparatus. 
\choice Actin and myosin filaments.  
\choice Basal apparatus and hook. 
\correctchoice Dynein and microtubules. 
\choice Sensory and motor neurons. 
\end{choices}

\newpage

\question
The fluid that moves around in the circulatory system of a typical arthropod is the

\begin{choices}
\choice cytoplasm. 
\choice blood plasma. 
\correctchoice hemolymph. 
\choice gastrovascular fluid. 
\choice intracellular fluid. 
\end{choices}

\question
Ignoring all other factors, what kind of day would result in the fastest delivery of water and minerals to the leaves of a tree?

\begin{choices}
\choice A cool, dry day. 
\choice A cool, humid day. 
\choice A very hot, dry, windy day. 
\correctchoice A warm, dry day. 
\choice A warm, humid day.  
\end{choices}

\question
The main stimulus used by birds to find their way north for the summer time is 

\begin{choices}
\choice thermal. 
\choice mechanical. 
\correctchoice electomagnetic.
\choice thigmic. 
\choice chemical. 
\end{choices}

\question
Most of the water taken up by a plant

\begin{choices}
\choice transports sugars in the phloem. 
\correctchoice is lost during transpiration. 
\choice is converted to sugars during photosynthesis. 
\choice enters the central vacuoles of cells.  
\choice keep guard cells turgid. 
\end{choices}

\question
The highest quality form of energy (best able to do work) for organisms is

\begin{choices}
\choice kinetic. 
\choice thermal. 
\correctchoice radiant. 
\choice chemical.
\choice potential.
\end{choices}

%\newpage
\question
Transpiration in plants requires all of the following \emph{except}

\begin{choices}
\choice adhesion of water molecules to cellulose. 
\choice evaporation of water molecules. 
\choice cohesion between water molecules.  
\correctchoice active transport through xylem cells. 
\choice water vapor on the spongy mesophyll. 
\end{choices}

\question
Organisms with a circulating fluid that is separate from the fluid that directly surrounding the body's cells are likely to have

\begin{choices}
\choice a gastrovascular cavity. 
\choice an open circulatory system. 
\choice a hemolymph system. 
\correctchoice a closed circulatory system. 
\choice a vascular system. 
\end{choices}



\question
Differences in energy quality and availability affects the amount of surface area and volume in organisms. Compared to animals, plants have 

\begin{choices}
\correctchoice high surface area and low volumes. 
\choice low surface area and high volumes. 
\choice low surface areas and low volumes.  
\choice high surface area and high volumes. 
\end{choices}


\question
What drives the flow of water through the xylem?

\begin{choices}
\correctchoice The evaporation of water from the leaves. 
\choice Abscisic acid opening guard cells around the stomata. 
\choice The number of companion cells in the phloem.  
\choice Transpirational pull through sieve-tube elements. 
\choice A higher concentration of sugars in the leaves than in the roots. 
\end{choices}

\question
The type of nerve cell responsible for integrating different stimuli is the 

\begin{choices}
\correctchoice interneuron 
\choice motor neuron 
\choice dendrite 
\choice sensory neuron 
\choice axon 
\end{choices}

\question
In terms of structure, cilia  are most like 
\begin{choices}
\choice microtubules.  
\choice bacterial flagella. 
\choice myosin filaments. 
\choice actin filaments. 
\correctchoice eukaryotic flagella. 
\end{choices}

\question
The ion used most often as a second messenger is

\begin{choices}
\choice \ce{K+} 
\choice \ce{H+} 
\correctchoice \ce{Ca+} 
\choice \ce{Cl-} 
\choice \ce{Na+} 
\end{choices}

\question
From an energy perspective, animals are mobile because they consume

\begin{choices}
\choice low quality energy and carbon intake is constant. 
\choice low quality energy and low energy efficiency.
\choice high quality energy and carbon intake is irregular. 
\choice high quality energy and carbon intake is constant. 
\correctchoice low quality energy and carbon intake is irregular.
\choice high quality energy and high energy efficiency. 
\end{choices}

\question
The function of valves in veins in a closed circulatory system is to 
\begin{choices}
\choice provide support for the extra muscles around the vein. 
\choice enhance diffusion of nutrients to the tissues. 
\choice decrease the amount of pressure in the veins.  
\choice enable the vein to withstand high pressure. 
\correctchoice ensure one way flow of circulatory fluids. 
\end{choices}

\question
The vascular tissue system for delivery of water, minerals, and nutrients throughout the plant includes all of the following except for

\begin{choices}
\choice xylem. 
\choice phloem. 
\choice sieve plates. 
\correctchoice cohesion. 
\choice companion cells. 
\end{choices}

\question
Hydrostatic skeletons are 

\begin{choices}
\choice skeletons made of microtubules in the cytoplasm of protists. 
\choice skeletons found in deep sea organisms that depend on high pressure of the surrounding water to provide support. 
\correctchoice fluid-filled compartments under pressure. 
\choice external skeletons that provide support and protection against dehydration. 
\choice internal skeletons that are buried in soft tissue.  
\end{choices}

\question
Galvanotaxis occurs when 

\begin{choices}
\choice plants respond to gravity.  
\choice plants respond to hormones. 
\choice animals respond to temperature. 
\choice animals respond to metal ions. 
\correctchoice protists respond to an electrical field. 
\choice protists respond to touch. 
\end{choices}

\newpage

\question
Pores on the leaf surface that function in gas and water exchange are called

\begin{choices}
\choice spongy mesophyll. 
\choice xylem. 
\choice phloem. 
\choice waxy cuticles. 
\correctchoice stomata. 
\end{choices}
